\documentclass[10pt,a4paper]{amsart}
\usepackage[margin=19mm]{geometry}
\usepackage{amsmath,amssymb,mathtools}
\usepackage[scr=rsfs,cal=euler]{mathalfa}
\usepackage{xfrac}
\usepackage[unicode,hidelinks,bookmarks=false]{hyperref}
\newcommand{\ie}{i.e.\ }
\newcommand{\cf}{cf.\ }
\newcommand{\resp}{respectively\ }
\newcommand{\Subsection}{\subsection}
\providecommand{\nobreakdash}{\nobreak\hbox{-}}
\setlength{\emergencystretch}{2em}
\multlinegap0pt
\usepackage{bm}
\usepackage{bbm}
\usepackage{dsfont}
\usepackage{xspace}
\usepackage{cancel}
\usepackage{mathtools}
\usepackage{amsthm}
\makeatletter
\@ifundefined{thm}{\newtheorem{thm}{Theorem}}{}
\@ifundefined{claim}{\newtheorem{claim}[thm]{Claim}}{}
\@ifundefined{lem}{\newtheorem{lem}[thm]{Lemma}}{}
\@ifundefined{rem}{\newtheorem{rem}[thm]{Remark}}{}
\makeatother
\newcommand{\Ric}{\mathrm{Ric} \,}
\newcommand{\f}{\varphi}
\newcommand{\p}{\psi}
\title[K\"ahler-Ricci flows coming out of metric spaces]{K\"ahler-Ricci flows coming out of metric spaces}
\author{Alix Deruelle, Vincent Guedj, Henri Guenancia, Ahmed Zeriahi}
\date{Source: arXiv:2511.13473v1 (CC BY 4.0)}
\begin{document}
\maketitle

\noindent\emph{Source and licence.} K\"ahler-Ricci flows coming out of metric spaces. Version arXiv:2511.13473v1 (CC BY 4.0), distributed under the Creative Commons Attribution 4.0 International licence, as recorded in the arXiv licence metadata for this exact version and shown on the abstract page https://arxiv.org/abs/2511.13473v1. Full rights notice: licenses/arxiv-2511.13473-CC-BY-4.0.txt.\par\smallskip

\noindent\textit{Editorial scope.} Counted unit: the Rem whose source label is \texttt{None} in the article (source0.tex lines 1334--1336, source label \texttt{None}), delivered together with the article's own complete original proof of it (source0.tex lines 1338--1445, one \texttt{proof} environment of 108 lines). No mathematics is written, repaired or re-derived: statement and proof are the article's own text line for line. Required context retained uncounted: lem:bddphidot (lines 955--960); thm:laplace (lines 1035--1040). Every definition, display and label of those lines is retained unchanged. Omitted from this package and not redistributed: 1--954, 961--1034, 1041--1333, 1337--1337, 1446--2714. Nothing inside a retained excerpt is omitted or altered: every retained body is delivered exactly as the article's lines are written, so no omission receipt of any kind accompanies this package. Citations: the retained excerpts carry no citation command at all, so no bibliography entry of the article is transcribed and no reference is redistributed. Reference adaptation: none was needed and none was made; every reference inside the delivered unit resolves inside it, so no unresolved reference or citation is produced. Printed slips kept literal: no printed word, symbol, hypothesis or number of the counted statement or of its proof was altered. Layout and class: the article's own class is \texttt{smfart} with packages including aeguill, enumerate, amssymb, amsmath, latexsym, amsthm, fontenc, geometry, url, babel, inputenc, mathrsfs, xcolor, comment; the wrapper is the lane's pinned \texttt{amsart} editorial wrapper at 10pt on A4, which restates the theorem-like environments the retained text needs and adds the article's own macro definitions that the retained text needs (\texttt{\textbackslash Ric}, \texttt{\textbackslash f}, \texttt{\textbackslash p}), with extra wrapper packages (bm, bbm, dsfont, xspace, cancel, mathtools). The delivered numbering is the unit's own. Assets: no retained span contains a figure environment, a graphics-inclusion command, a PDF-page inclusion, a table or a TikZ picture; no image file is copied into this package, the package declares no asset dependency and no image file is redistributed with it. Licence: version arXiv:2511.13473v1 of this article is distributed under the Creative Commons Attribution 4.0 International licence, as recorded in the arXiv licence metadata for this exact version and shown on the abstract page https://arxiv.org/abs/2511.13473v1. A full rights notice is delivered at licenses/arxiv-2511.13473-CC-BY-4.0.txt. Characters: every retained excerpt is pure ASCII, so the delivered engine, XeLaTeX with its default Latin Modern text font, reports no missing character.
\begin{lem} \label{lem:bddphidot}
There exists $C>0$ such that for all $0 \leq t \leq 1$ and for all $x \in X$,
$$
 {\p_+}(x) -C   \leq   \dot{\f_t}(x) \leq C-{\p_-}(x) .
$$
\end{lem}

\begin{thm} \label{thm:laplace}
There exists $C>0$ such that for all $0 \leq t \leq 1$ and for all $x \in X$,
$$
\frac{1}{C} e^{\p_+} \omega_X \leq \omega_X+dd^c \f_t \leq C e^{-\p_-} \omega_X.
$$
\end{thm}

\begin{rem}
Note that the gradient of $\dot{\f_t}$ does not blow up in time as $t$ tends to $0^+$.
\end{rem}

\begin{proof}
Denote by $g(t)$ and $g_X$ respectively the Riemannian metrics induced by the K\"ahler forms $\omega(t)$ and $\omega_X$.
For $\gamma_0>0$, compute as follows:
\begin{equation}
\begin{split}\label{it-never-ends}
\left(\partial_t-\frac{1}{2}\Delta_{g(t)}\right)\left(e^{\gamma_0\psi_{-}}\dot{\f_t}^2\right)&=\left(-\frac{1}{2}\Delta_{g(t)}e^{\gamma_0\psi_{-}}\right)\dot{\f_t}^2-\gamma_0g(t)\left(\nabla^{g(t)}\psi_{-},\nabla^{g(t)}\dot{\f_t}^2\right)e^{\gamma_0\psi_{-}}\\
&\quad-\,e^{\gamma_0\psi_{-}}|\nabla^{g(t)}\dot{\f_t}|^2_{g(t)}.
\end{split}
\end{equation}
Now, as $i\partial\overline{\partial}\psi_{-}+A\omega_X\geq 0$ for some $A> 0$, 
\begin{equation*}
\begin{split}
-\Delta_{g(t)}e^{\gamma_0\psi_{-}}&=\left(-\gamma_0\Delta_{g(t)} \psi_{-}-\gamma_0^2|\nabla^{g(t)}\psi_{-}|^2_{g(t)}\right)e^{\gamma_0\psi_{-}}\leq C\gamma_0 Ae^{\gamma_0\psi_{-}},
\end{split}
\end{equation*}
where $C$ is a positive constant that may vary from line to line but which is independent of $t>0$. Here we have invoked Theorem \ref{thm:laplace} in the last inequality. 

Therefore, Young's inequality applied to the mixed term of \eqref{it-never-ends} gives us:
\begin{equation}
\begin{split}\label{est-0-undergrad}
\left(\partial_t-\frac{1}{2}\Delta_{g(t)}\right)\left(e^{\gamma_0\psi_{-}}\dot{\f_t}^2\right)&
\leq C\left(\gamma_0A+\gamma_0^2|\nabla^{g(t)}\psi_{-}|^2_{g(t)}\right)e^{\gamma_0\psi_{-}}\dot{\f_t}^2 -\frac{1}{2}e^{\gamma_0\psi_{-}}|\nabla^{g(t)}\dot{\f_t}|^2_{g(t)}\\
&\leq C -\frac{1}{2}e^{\gamma_0\psi_{-}}|\nabla^{g(t)}\dot{\f_t}|^2_{g(t)} ,
\end{split}
\end{equation}
where $C$ may vary from line to line. Here we have used  $|\nabla^{g(t)}\psi_{-}|_{g(t)}\leq C|\nabla^{g_X}\psi_{-}|_{g_X}$ according to Theorem \ref{thm:laplace} and the fact that $\psi_+=0$ in the last line: indeed  $\dot{\f_t}^2\leq C(\psi_{-})^2$ according to Lemma \ref{lem:bddphidot} together with the geometric assumption on the $g_X$-gradient of $\psi_{-}$, one concludes as above by choosing $\gamma_0>0$ large enough compared to $\alpha$. 

\begin{claim}\label{claim-bochner}
There exists $\gamma_0>0$ such that for $\gamma_1>\gamma_0$, there exists $C>0$ such that,
\begin{equation}
\begin{split}\label{est-I-grad}
\left(\partial_t-\frac{1}{2}\Delta_{g(t)}\right)\left(e^{\gamma_1\psi_{-}}|\nabla^{g(t)}\dot{\f_t}|_{g(t)}^2\right)&\leq Ce^{\gamma_0\psi_{-}}|\nabla^{g(t)}\dot{\f_t}|_{g(t)}^2.
\end{split}
\end{equation}
\end{claim}

\begin{proof}[Proof of Claim \ref{claim-bochner}]
Observe that by a straightforward computation:
{\small
\begin{equation*}
\begin{split}
&\left(\partial_t-\frac{1}{2}\Delta_{g(t)}\right)\left(e^{\gamma_1\psi_{-}}|\nabla^{g(t)}\dot{\f_t}|_{g(t)}^2\right)\\
&=\left(\left(\partial_t-\frac{1}{2}\Delta_{g(t)}\right)e^{\gamma_1\psi_{-}}\right)|\nabla^{g(t)}\dot{\f_t}|_{g(t)}^2-g(t)\left(\nabla^{g(t)}e^{\gamma_1\psi_{-}},\nabla^{g(t)}|\nabla^{g(t)}\dot{\f_t}|_{g(t)}^2\right)\\
&\quad+e^{\gamma_1\psi_{-}}\left(\partial_t-\frac{1}{2}\Delta_{g(t)}\right) |\nabla^{g(t)}\dot{\f_t}|_{g(t)}^2\\
&=\left(-\frac{\gamma_1}{2}\Delta_{g(t)}\psi_{-}-\gamma_1^2|\nabla^{g(t)}\psi_{-}|^2_{g(t)}\right)e^{\gamma_1\psi_{-}}|\nabla^{g(t)}\dot{\f_t}|_{g(t)}^2\\
&\quad-\gamma_1e^{\gamma_1\psi_{-}}g(t)\left(\nabla^{g(t)}\psi_{-},\nabla^{g(t)}|\nabla^{g(t)}\dot{\f_t}|_{g(t)}^2\right)+e^{\gamma_1\psi_{-}}\left(\partial_t-\frac{1}{2}\Delta_{g(t)}\right) |\nabla^{g(t)}\dot{\f_t}|_{g(t)}^2.
\end{split}
\end{equation*}
}
Therefore, Young's inequality implies,
{\small
\begin{equation}
\begin{split}\label{lovely-formula}
&\left(\partial_t-\frac{1}{2}\Delta_{g(t)}\right)\left(e^{\gamma_1\psi_{-}}|\nabla^{g(t)}\dot{\f_t}|_{g(t)}^2\right)\leq\left(C\gamma_1 A-\frac{\gamma_1^2}{2}|\nabla^{g(t)}\psi_{-}|^2_{g(t)}\right)e^{\gamma_1\psi_{-}}|\nabla^{g(t)}\dot{\f_t}|_{g(t)}^2\\
&\quad+2\gamma_1e^{\gamma_1\psi_{-}}|\nabla^{g(t)}\psi_{-}|_{g(t)}|\nabla^{g(t)}\dot{\f_t}|_{g(t)}|\nabla^{g(t),\,2}\dot{\f_t}|_{g(t)}+e^{\gamma_1\psi_{-}}\left(\partial_t-\frac{1}{2}\Delta_{g(t)}\right) |\nabla^{g(t)}\dot{\f_t}|_{g(t)}^2\\
&\leq C\left(A+|\nabla^{g(t)}\psi_{-}|^2_{g(t)}\right)e^{\gamma_1\psi_{-}}|\nabla^{g(t)}\dot{\f_t}|_{g(t)}^2\\
&\quad +e^{\gamma_1\psi_{-}}|\nabla^{g(t),\,2}\dot{\f_t}|^2_{g(t)}+e^{\gamma_1\psi_{-}}\left(\partial_t-\frac{1}{2}\Delta_{g(t)}\right) |\nabla^{g(t)}\dot{\f_t}|_{g(t)}^2,
\end{split}
\end{equation}
}
where $C$ is a positive constant depending on $\gamma_1$ that may vary from line to line.

Now, the Bochner formula together with the fact that $\partial_tg(t)=-\Ric(g(t))+\Ric(g_X)$ allow us to write:
\begin{equation*}
\begin{split}
\left(\partial_t-\frac{1}{2}\Delta_{g(t)}\right)|\nabla^{g(t)}\dot{\f_t}|^2_{g(t)}&=-|\nabla^{g(t),\,2}\dot{\f_t}|^2_{g(t)}-\Ric(g_X)(\nabla^{g(t)}\dot{\f_t},\nabla^{g(t)}\dot{\f_t})\\
&\leq -|\nabla^{g(t),\,2}\dot{\f_t}|^2_{g(t)}+C|\nabla^{g(t)}\dot{\f_t}|^2_{g_X}\\
&\leq -|\nabla^{g(t),\,2}\dot{\f_t}|^2_{g(t)}+C|\nabla^{g(t)}\dot{\f_t}|^2_{g(t)},
\end{split}
\end{equation*}
where $C$ is a positive constant that may vary from line to line. Here we have used Theorem \ref{thm:laplace} in the last line.

Finally, the combination of \eqref{lovely-formula} and the previous estimate gives: 
\begin{equation*}
\begin{split}
\left(\partial_t-\frac{1}{2}\Delta_{g(t)}\right)\left(e^{\gamma_1\psi_{-}}|\nabla^{g(t)}\dot{\f_t}|_{g(t)}^2\right)&\leq C\left(A+|\nabla^{g(t)}\psi_{-}|^2_{g(t)}\right)e^{\gamma_1\psi_{-}}|\nabla^{g(t)}\dot{\f_t}|_{g(t)}^2\\
&\leq Ce^{\gamma_0\psi_{-}}|\nabla^{g(t)}\dot{\f_t}|_{g(t)}^2,
\end{split}
\end{equation*}
since by Theorem \ref{thm:laplace} together with the geometric assumption on the $g_X$-gradient of $\psi_{-}$, 
$$
|\nabla^{g(t)}\psi_{-}|^2_{g(t)}e^{(\gamma_1-\gamma_0)\psi_{-}}
\leq C|\nabla^{g_X}\psi_{-}|_{g_X}^2e^{(\gamma_1-\gamma_0)\psi_{-}}\leq C.
$$
%where $C$ is a positive constant that may vary from an inequality to another. 
This ends the proof of the claim.
%Notice here that we have dropped the good non-positive term containing $|\nabla^{g(t),\,2}\dot{\f_t}|_{g(t)}^2$.
\end{proof}

Let us sum the estimates  \eqref{est-0-undergrad} and [\eqref{est-I-grad}, Claim \ref{claim-bochner}] to get for $B>0$ large enough:
\begin{equation*}
\begin{split}\label{est-I-grad-bis}
&\left(\partial_t-\frac{1}{2}\Delta_{g(t)}\right)\left(B\,e^{\gamma_0\psi_{-}}\dot{\f_t}^2+e^{\gamma_1\psi_{-}}|\nabla^{g(t)}\dot{\f_t}|_{g(t)}^2\right)\\
&\leq BC -\frac{B}{2}e^{\gamma_0\psi_{-}}|\nabla^{g(t)}\dot{\f_t}|^2_{g(t)} +Ce^{\gamma_0\psi_{-}}|\nabla^{g(t)}\dot{\f_t}|_{g(t)}^2
\leq BC.
\end{split}
\end{equation*}

 Assuming $\gamma_1>0$ is chosen such that  
% $$
\begin{eqnarray*}
 \limsup_{t\rightarrow 0^+}e^{\gamma_1\psi_{-}}|\nabla^{g(t)}\dot{\f_t}|_{g(t)}^2
 &\leq & \limsup_{t\rightarrow 0^+}e^{\gamma_1\psi_{-}}|\nabla^{g_X}\dot{\f_t}|_{g_X}^2 \\
 &=& e^{\gamma_1\psi_{-}}|\nabla^{g_X}\psi_{-}|_{g_X}^2\leq C,
 \end{eqnarray*}
% $$ 
 which is made possible by assumption on $\psi_{-}$, the maximum principle applied to $$B\,e^{\gamma_0\psi_{-}}\dot{\f_t}^2+e^{\gamma_1\psi_{-}}|\nabla^{g(t)}\dot{\f_t}|_{g(t)}^2-Be^{Ct},$$ gives the desired result.
\end{proof}

\end{document}
